\section{Introduction}

The most difficult errors to repair are precisely the ones that provide the least actionable feedback. Assertion errors constitute 53\% of failures on standard code generation benchmarks and, according to prior work, exhibit among the lowest repair success rates, yet their error messages contain near-zero extractable state information. This paradox---that the errors most in need of rich diagnostic feedback are those that provide the least---reveals a fundamental tension in LLM-based code repair.

Recent advances in self-repair demonstrate that iterative refinement with verification feedback can improve pass rates by 4.9--17.1 percentage points on HumanEval \citep{howmanytries2026}. Runtime traces outperform simple error messages \citep{debugrepair2026}, and type constraints improve functional correctness \citep{tyflow2025}. However, these findings describe \emph{which} feedback works better without explaining \emph{why}. What properties of a verification signal enable an LLM to localize bugs and infer correct fixes?

We hypothesize that \textbf{Actionable Specificity (AS)} mediates feedback effectiveness. We decompose AS into three measurable components: localization (AS\_loc)---whether file:line information is present; state exposure (AS\_state)---how many variable values are visible; and causal context (AS\_causal)---execution trace depth between assignment and failure. Under this framework, a full runtime trace (high AS) should outperform a bare assertion error (low AS) because it provides the LLM with localization, state, and causal information necessary for targeted repair.

To test whether AS is even measurable from standard verification signals---a necessary precondition for testing its predictive power---we designed an existence test on HumanEval and MBPP (664 problems). We generated 600 signals across 6 conditions (C1--C6) varying in expected AS level, then extracted AS components using regex patterns designed for Python tracebacks.

\textbf{Our key finding is negative:} regex-based extraction achieves only 49.5\% coverage, with AS\_state at 2.7\%. The root cause is that assertion errors, which dominate these benchmarks (53.4\%), produce minimal output without traceback frames. Standard \texttt{AssertionError} messages lack the file:line and variable context that our extraction patterns require.

\textbf{We acknowledge this as a limitation of our operationalization.} However, the failure is informative: it reveals a boundary condition that future work must address. We demonstrate that:

\begin{enumerate}
    \item \textbf{Error type stratification is necessary}: feedback informativeness research must verify extraction feasibility by error type before operationalizing.
    \item \textbf{AS ordering is preserved} (C1$\geq$C2$\geq$C3$\geq$C4) on the extractable subset, suggesting the conceptual decomposition may remain valid under alternative operationalizations.
    \item \textbf{The AS framework provides actionable guidance}: it identifies \emph{which} error types block extraction and \emph{why} (missing file:line structure), enabling targeted solutions (e.g., \texttt{pytest --tb=long}, LLM-based extraction).
\end{enumerate}

The contribution is methodological: we document a necessary precondition---extraction feasibility by error type---that researchers must verify before designing feedback mechanisms for LLM code repair.
